% from ega2-1-fr.tex
\oldpage[II]{62}

affine $\Phi:G(\varphi)\to X'$, qui est dit \emph{associé} à $\varphi$ et que
nous noterons $\operatorname{Proj}(\varphi)$. Lorsque, pour tout $y\in Y$, il
y a un voisinage affine $U$ de $y$ tel que le
$\Gamma(U,\mathcal{O}_Y)$-module $\Gamma(U,\mathcal{S}_+)$ soit engendré par
$\varphi(\Gamma(U,\mathcal{S}'_+))$, on a
$G(\varphi_U)=p^{-1}(U)$, et par suite $G(\varphi)=X$.
\end{env}

\begin{proposition}[3.5.2]
\label{II.3.5.2-fr}
\begin{enumerate}
  \item[\rm(i)] Si $\mathcal{M}$ est un $\mathcal{S}$-Module gradué
    quasi-cohérent, il existe un isomorphisme canonique fonctoriel du
    $\mathcal{O}_{X'}$-Module $(\mathcal{M}_{[\varphi]})^{\sim}$ sur le
    $\mathcal{O}_{X'}$-Module
    $\Phi_*(\widetilde{\mathcal{M}}|G(\varphi))$.
  \item[\rm(ii)] Si $\mathcal{M}'$ est un $\mathcal{S}'$-Module gradué
    quasi-cohérent, il existe un homomorphisme canonique fonctoriel $\nu$ du
    $(\mathcal{O}_X|G(\varphi))$-Module
    $\Phi^*(\widetilde{\mathcal{M}'})$ dans le
    $(\mathcal{O}_X|G(\varphi))$-Module
    $(\mathcal{M}'\otimes_{\mathcal{S}'}\mathcal{S})^{\sim}|G(\varphi)$. Si
    $\mathcal{S}'$ est engendrée par $\mathcal{S}'_1$, $\nu$ est un
    isomorphisme.
\end{enumerate}
\end{proposition}

\begin{proof}
Les homomorphismes considérés ont en effet déjà été définis lorsque $Y$ est
affine (\hyperref[II.2.8.7-fr]{2.8.7} et
\hyperref[II.2.8.8-fr]{2.8.8}), et dans le cas général il suffit de vérifier
qu'ils sont compatibles avec les opérations de restriction d'un ouvert
affine de $Y$ à un ouvert plus petit, ce qui résulte aussitôt de la
commutativité de (\hyperref[II.3.5.1.1-fr]{3.5.1.1}).
\end{proof}

En particulier, pour tout $n\in\mathbf{Z}$, on a un homomorphisme canonique
\[
\label{II.3.5.2.1-fr}
  \Phi^*(\mathcal{O}_{X'}(n))\longrightarrow
  \mathcal{O}_X(n)|G(\varphi).
\tag{3.5.2.1}
\]

\begin{proposition}[3.5.3]
\label{II.3.5.3-fr}
Soient $Y$, $Y'$ deux préschémas, $\psi:Y'\to Y$ un morphisme,
$\mathcal{S}$ une $\mathcal{O}_Y$-Algèbre graduée quasi-cohérente, et posons
$\mathcal{S}'=\psi^*(\mathcal{S})$. Alors le $Y'$-schéma
$X'=\operatorname{Proj}(\mathcal{S}')$ s'identifie canoniquement à
$\operatorname{Proj}(\mathcal{S})\times_Y Y'$. En outre, si $\mathcal{M}$ est
un $\mathcal{S}$-Module gradué quasi-cohérent, le $\mathcal{O}_{X'}$-Module
$(\psi^*(\mathcal{M}))^{\sim}$ s'identifie à
$\widetilde{\mathcal{M}}\otimes_Y\mathcal{O}_{Y'}$.
\end{proposition}

\begin{proof}
Notons en premier lieu que $\psi^*(\mathcal{S})$ et
$\psi^*(\mathcal{M})$ sont des $\mathcal{O}_{Y'}$-Modules quasi-cohérents
ainsi que leurs composants homogènes
(\hyperref[0.5.1.4-fr]{0, 5.1.4}). Soient $U$ un ouvert affine de $Y$,
$U'\subset\psi^{-1}(U)$ un ouvert affine de $Y'$, $A$, $A'$ les anneaux de
$U$ et $U'$ respectivement~; on a alors
$\mathcal{S}|U=\widetilde{S}$, où $S$ est une $A$-algèbre graduée, et
$\mathcal{S}'|U'$ s'identifie à $(S\otimes_A A')^{\sim}$
(\hyperref[I.1.6.5-fr]{I, 1.6.5})~; la première assertion résulte alors de
(\hyperref[II.2.8.10-fr]{2.8.10}) et de
(\hyperref[I.3.2.6.2-fr]{I, 3.2.6.2}), car on vérifie aussitôt que la
projection
$\operatorname{Proj}(\mathcal{S}'|U')\to
\operatorname{Proj}(\mathcal{S}|U)$ définie par l'identification précédente
est compatible avec les opérations de restrictions sur $U$ et $U'$ et
définit donc bien un morphisme
$\operatorname{Proj}(\mathcal{S}')\to\operatorname{Proj}(\mathcal{S})$.
Soient maintenant
\[
  q:\operatorname{Proj}(\mathcal{S})\to Y,
  \qquad q':\operatorname{Proj}(\mathcal{S}')\to Y'
\]
les morphismes structuraux~; $q'^{-1}(U')$ s'identifie alors à
$q^{-1}(U)\times_U U'$, et les deux faisceaux
$(\psi^*(\mathcal{M}))^{\sim}|q'^{-1}(U')$ et
$(\widetilde{\mathcal{M}}\otimes_Y\mathcal{O}_{Y'})|q'^{-1}(U')$
s'identifient alors canoniquement tous deux à $(M\otimes_A A')^{\sim}$, où
on a posé $M=\Gamma(U,\mathcal{M})$, en vertu de
(\hyperref[II.2.8.10-fr]{2.8.10}) et de
(\hyperref[I.1.6.5-fr]{I, 1.6.5})~; d'où la seconde assertion, car on vérifie
encore de façon immédiate la compatibilité des identifications précédentes
avec les opérations de restriction.
\end{proof}

\begin{corollary}[3.5.4]
\label{II.3.5.4-fr}
Avec les notations de (\hyperref[II.3.5.3-fr]{3.5.3}),
$\mathcal{O}_{X'}(n)$ s'identifie canoniquement à
$\mathcal{O}_X(n)\otimes_Y\mathcal{O}_{Y'}$ pour tout $n\in\mathbf{Z}$ (avec
$X=\operatorname{Proj}(\mathcal{S})$).
\end{corollary}

\begin{proof}
En effet, avec les notations de (\hyperref[II.3.5.3-fr]{3.5.3}), il est clair
que $\psi^*(\mathcal{S}(n))=\mathcal{S}'(n)$ pour tout $n\in\mathbf{Z}$.
\end{proof}

\begin{env}[3.5.5]
\label{II.3.5.5-fr}
Gardant les notations précédentes, désignons par $\Psi$ la projection
canonique $X'\to X$, et posons $\mathcal{M}'=\psi^*(\mathcal{M})$~; nous
supposerons en outre que $\mathcal{S}$ est engendrée par $\mathcal{S}_1$

